
% ============================
% CHAPTERS 90–95 FOR THE TOME
% Part XIX — Autonomy, Free-Will Geometry, and Reflective Cognitive Fields
% ============================

\part{Autonomy, Free-Will Geometry, and Reflective Cognitive Fields}

% ---------- CHAPTER 90 ----------
\chapter{Autonomy as a Phase-Separated Cognitive Field}

\section{Autonomy Field Definition}
Define an autonomy field
\[
\tau_{\mathrm{aut}}(x,t) = \tau_{\mathrm{cog}}(x,t) - \tau_{\mathrm{ext}}(x,t),
\]
where \(\tau_{\mathrm{cog}}\) encodes internal cognitive drives and \(\tau_{\mathrm{ext}}\) encodes external constraint fields.

\section{Autonomy Magnitude}
The local autonomy magnitude is
\[
A_{\mathrm{loc}}(x,t) = (\partial\tau_{\mathrm{aut}})^2,
\]
and the global autonomy functional
\[
\mathcal{A} = \int A_{\mathrm{loc}}(x,t)\, dV
\]
quantifies the degree of internally driven behaviour.

% ---------- CHAPTER 91 ----------
\chapter{Free-Will Geometry and Competing Phase Potentials}

\section{Competing Potentials}
Let \(V_{\mathrm{int}}(\tau)\) encode internal preference landscapes and \(V_{\mathrm{ext}}(\tau)\) encode external pressures.
The effective decision potential is
\[
V_{\mathrm{eff}}(\tau) = V_{\mathrm{int}}(\tau) + V_{\mathrm{ext}}(\tau).
\]

\section{Decision Trajectories}
A decision variable \(q(t)\) evolves according to
\[
\partial_t q = -\mu\,\partial_q V_{\mathrm{eff}}(q),
\]
with “free” decisions corresponding to regimes where \(|\partial_q V_{\mathrm{int}}| \gg |\partial_q V_{\mathrm{ext}}|\).

% ---------- CHAPTER 92 ----------
\chapter{Reflective Fields and Meta-Cognitive Feedback}

\section{Reflective Phase Component}
Define a reflective component
\[
\tau_{\mathrm{refl}}(x,t) = \mathcal{R}[\tau_{\mathrm{cog}}](x,t),
\]
where \(\mathcal{R}\) is a feedback operator acting on the cognitive field.

\section{Feedback Dynamics}
A simple feedback law is
\[
\partial_t \tau_{\mathrm{refl}} = -\gamma\left(\tau_{\mathrm{refl}} - \tau_{\mathrm{cog}}\right) + \eta\,\partial_t \tau_{\mathrm{cog}},
\]
which smooths and tracks cognitive evolution, enabling self-monitoring and error correction.

% ---------- CHAPTER 93 ----------
\chapter{Constraint Fields, Agency, and Environmental Coupling}

\section{Constraint Field}
External constraints (social, physical, institutional) are encoded in
\[
\tau_{\mathrm{con}}(x,t),
\]
with associated constraint energy
\[
E_{\mathrm{con}} = \int (\partial\tau_{\mathrm{con}})^2 dV.
\]

\section{Agency Measure}
Agency is defined as
\[
\mathcal{G} = \int (\partial\tau_{\mathrm{cog}}\cdot \partial\tau_{\mathrm{env}})\, dV,
\]
where \(\tau_{\mathrm{env}}\) describes the environment.
Positive \(\mathcal{G}\) indicates alignment between internal drives and environmental affordances; negative \(\mathcal{G}\) signals conflict.

% ---------- CHAPTER 94 ----------
\chapter{Phase-Compatible Environments and Supportive Cognitive Geometry}

\section{Environmental Matching}
A phase-compatible environment satisfies
\[
(\partial\tau_{\mathrm{cog}} - \partial\tau_{\mathrm{env}})^2 < \epsilon_{\mathrm{match}},
\]
reducing cognitive strain and enhancing stability.

\section{Supportive Design Principle}
Supportive environments minimise the functional
\[
\mathcal{S}_{\mathrm{env}} = \int (\partial\tau_{\mathrm{cog}} - \partial\tau_{\mathrm{env}})^2 dV,
\]
providing a geometric criterion for designing architectures, sensory contexts, and interaction patterns that reduce overload and distress.

% ---------- CHAPTER 95 ----------
\chapter{Adaptive Phase Interfaces and Assistive Cognitive Technologies}

\section{Adaptive Interface Field}
An assistive interface couples to the cognitive field via
\[
\tau_{\mathrm{int}}(x,t) = \alpha\,\tau_{\mathrm{cog}}(x,t) + \beta\,\tau_{\mathrm{env}}(x,t),
\]
and adjusts its output to minimise
\[
\mathcal{E}_{\mathrm{int}} = \int (\partial\tau_{\mathrm{int}} - \partial\tau_{\mathrm{target}})^2 dV,
\]
where \(\tau_{\mathrm{target}}\) encodes a desired comfort or focus state.

\section{Toward Neuro-Supportive Devices}
By continuously estimating \(\tau_{\mathrm{cog}}\) and \(\tau_{\mathrm{env}}\), adaptive phase interfaces can in principle:
\begin{itemize}
  \item damp disruptive gradients (e.g.\ sensory overload),
  \item enhance stabilising modes (e.g.\ focus, calm),
  \item and provide a real-time, geometry-based feedback channel between mind and environment.
\end{itemize}
This chapter lays the groundwork for future MK--43 Ultra inspired assistive technologies, including advanced auditory and multisensory modulation systems.
